\chapter{Comida y bebida}
\section{Objetivos}
\begin{itemize}[leftmargin=*]
	\item Nombrar alimentos y bebidas comunes y hacer pedidos simples.
	\item Practicar cantidades básicas y preferencias.
\end{itemize}

\section{Contenidos clave}
\begin{itemize}[leftmargin=*]
	\item Vocabulario: comidas del día, supermercado, café/restaurante.
	\item Estructuras: \textit{Ik wil graag \ldots{}}, \textit{Mag ik \ldots{}?}, \textit{Ik houd van \ldots{} / Ik lust geen \ldots{}}.
	\item Cantidades: \textit{een glas, een kop, een broodje, kilo, liter}.
\end{itemize}

\section{Ruta del capítulo}
Seis secciones con vocabulario, frases modelo y práctica aplicada.

\section{Comidas del día}
\textbf{Términos básicos.}
\begin{itemize}[leftmargin=*]
  \item \textit{ontbijt} (desayuno), \textit{lunch} / \textit{middageten} (almuerzo), \textit{avondeten} (cena), \textit{tussendoortje} (snack).
  \item Verbos: \textit{ontbijten} (desayunar), \textit{lunchen}, \textit{avondeten} / \textit{dineren} (formal).
\end{itemize}

\textbf{Frases modelo}
\begin{itemize}[leftmargin=*]
  \item \textit{Ik ontbijt om zeven uur}. (Desayuno a las 7.)
  \item \textit{Wij lunchen op kantoor}. (Almorzamos en la oficina.)
  \item \textit{Wat eet je vanavond?} (¿Qué comes esta noche?)
\end{itemize}

\textbf{Mini práctica}
\begin{itemize}[leftmargin=*]
  \item Escribe tu horario de comidas con horas en formato 24h.
  \item Pregunta a un compañero qué suele comer en cada comida.
\end{itemize}

\section{Alimentos básicos}
\textbf{Vocabulario A1.}
\begin{itemize}[leftmargin=*]
  \item Pan y lácteos: \textit{brood, boter, kaas, melk, yoghurt}.
  \item Frutas y verduras: \textit{appel, banaan, sinaasappel, aardappel, tomaat, sla, ui}.
  \item Proteínas: \textit{kip, vis, ei, vlees, kaas}.
  \item Dulces y snacks: \textit{koekje, taart, chips, chocolade}.
\end{itemize}

\textbf{Estructuras}
\begin{itemize}[leftmargin=*]
  \item \textit{Ik eet graag \\underline{\hspace{1.2cm}}}. (Me gusta comer \underline{\hspace{1.2cm}}.)
  \item \textit{Ik koop brood en kaas}. (Compro pan y queso.)
\end{itemize}

\textbf{Mini práctica}
\begin{itemize}[leftmargin=*]
  \item Escribe tu lista de 8 alimentos básicos en neerlandés.
  \item Forma dos frases con \textit{graag} para preferencias.
\end{itemize}

\section{Bebidas comunes}
\textbf{Vocabulario.}
\begin{itemize}[leftmargin=*]
  \item Calientes: \textit{koffie, thee, warme chocolademelk}.
  \item Frías: \textit{water, sap, frisdrank, bier, wijn}.
  \item Pedir tamaño: \textit{een kleine / grote koffie}, \textit{een glas water}, \textit{een flesje cola}.
\end{itemize}

\textbf{Frases modelo}
\begin{itemize}[leftmargin=*]
  \item \textit{Mag ik een koffie zonder suiker?}
  \item \textit{Ik drink geen alcohol}. (No bebo alcohol.)
  \item \textit{Wil je nog iets drinken?}
\end{itemize}

\textbf{Mini práctica}
\begin{itemize}[leftmargin=*]
  \item Pide dos bebidas distintas con cortesía \textit{alstublieft/alsjeblieft}.
  \item Niega una bebida que no tomas habitualmente.
\end{itemize}

\section{Pedir en un restaurante}
\textbf{Pasos y frases útiles.}
\begin{itemize}[leftmargin=*]
  \item Llamar al camarero: \textit{Ober! Mag ik de menukaart, alstublieft?}
  \item Pedir: \textit{Ik wil graag de dagsoep} / \textit{Mag ik een broodje kaas?}
  \item Confirmar o cambiar: \textit{Heeft u ook vegetarisch?} \textit{Zonder ui, alstublieft}.
  \item La cuenta: \textit{Mag ik de rekening, alstublieft?}
\end{itemize}

\textbf{Propinas y pago}
\begin{itemize}[leftmargin=*]
  \item Propina no obligatoria; se suele redondear.~\textit{Mag ik met pin betalen?}
\end{itemize}

\textbf{Mini práctica}
\begin{itemize}[leftmargin=*]
  \item Role-play de pedido con entrada, plato y bebida.
  \item Escribe tres preguntas al camarero sobre el menú.
\end{itemize}

\section{Expresar gustos y preferencias}
\textbf{Verbos y partículas.}
\begin{itemize}[leftmargin=*]
  \item \textit{Ik houd van \\_} (me gusta), \textit{Ik hou niet van \\_} (no me gusta).
  \item \textit{Ik lust \\_ (niet)} se usa con comida/bebida que (no) te gusta consumir.
  \item Intensificadores: \textit{graag, heel graag, helemaal niet}.
\end{itemize}

\textbf{Ejemplos}
\begin{itemize}[leftmargin=*]
  \item \textit{Ik houd van kaas, maar ik hou niet van vis}. 
  \item \textit{Ik lust geen koffie}. (No me gusta el café.)
  \item \textit{Wij eten graag pasta}. 
\end{itemize}

\textbf{Mini práctica}
\begin{itemize}[leftmargin=*]
  \item Escribe tres frases: un gusto fuerte, un disgusto y algo neutral.
  \item Cambia \textit{houden van} a \textit{lusten} cuando hables de alimentos/bebidas.
\end{itemize}

\section{Compras en el supermercado}
\textbf{Vocabulario práctico.}
\begin{itemize}[leftmargin=*]
  \item Secciones: \textit{groente, fruit, vlees, zuivel, brood, diepvries}.
  \item Envases y cantidades: \textit{een pak melk, een kilo appels, een pond kaas, een blik tomaten, een zak chips}.
  \item Ofertas: \textit{in de aanbieding}, 2 voor 1: \textit{twee voor de prijs van één}.
\end{itemize}

\textbf{Frases modelo}
\begin{itemize}[leftmargin=*]
  \item \textit{Waar ligt de melk?} (¿Dónde está la leche?)
  \item \textit{Ik neem een kilo tomaten en twee pakken pasta}. 
  \item \textit{Mag ik met pin betalen?} (¿Puedo pagar con tarjeta?)
\end{itemize}

\textbf{Mini práctica}
\begin{itemize}[leftmargin=*]
  \item Crea una lista de compra de 8 ítems con cantidades.
  \item Role-play cajero/cliente: pagar y pedir bolsa.
\end{itemize}


\section{Práctica sugerida}
\begin{itemize}[leftmargin=*]
	\item Simulación de pedir en cafetería y leer un menú simple.
	\item Lista de compras con precios estimados.
	\item Recetas cortas con imperativos básicos (exposición).
\end{itemize}

\section{Microevaluación}
\begin{itemize}[leftmargin=*]
	\item Role-play de pedido en 6 líneas; lista de 10 ítems de compra.
\end{itemize}
